\paragrafos{Βασικά στερεά}{Β.4.1, «Βασικά στερεά: ορθό πρίσμα, παραλληλεπίπεδο, κύβος, πυραμίδα, κύλινδρος, κώνος, σφαίρα»}

\begin{anagke}
Η Αγγελική θέλει να φτιάξει ένα κουτί δώρου για τα γενέθλια του παππού Γιάννη. Παίρνει χαρτόνι, κόβει έξι τετράγωνα και τα βάζει στη σειρά, το ένα δίπλα στο άλλο. Όταν προσπαθεί να τα διπλώσει, τα τετράγωνα πέφτουν το ένα πάνω στο άλλο και το κουτί μένει ανοιχτό. «Δεν κλείνει!» λέει απογοητευμένη. Ο Ηλίας κόβει άλλα έξι τετράγωνα και τα βάζει σε σχήμα σταυρού. Τώρα διπλώνονται ακριβώς σε κύβο.

Στη βιτρίνα του μαγαζιού απέναντι η Ρήνη βλέπει κουτιά παπουτσιών, τενεκεδάκια με φασόλια, καπέλα για πάρτι και μια μπάλα. Ο Στέλιος λέει με σιγουριά: «Όλα αυτά είναι κουτιά». Η κυρία Αλεξάνδρα χαμογελά: «Για δες καλύτερα. Πόσες έδρες έχει το καθένα; Έχει γωνίες; Μπορείς να το ανοίξεις και να το απλώσεις στο τραπέζι;». Η Πολυξένη γράφει στο τετράδιό της έναν πίνακα με τα ονόματα και τα στοιχεία κάθε σχήματος.

Στη γεωμετρία του χώρου δεν δουλεύουμε μόνο με σχήματα που ζωγραφίζουμε σε χαρτί, αλλά και με σώματα που έχουν μήκος, πλάτος και ύψος. Σε αυτή την παράγραφο θα γνωρίσουμε τα βασικά στερεά, τα στοιχεία τους, τα αναπτύγματά τους και τον τρόπο που τα σχεδιάζουμε.
\end{anagke}

\begin{lexi}{στερεό, πρίσμα, κώνος, ανάπτυγμα}{ονόματα στερεών}
\lexiword{Στερεό} λέμε ένα σώμα που καταλαμβάνει χώρο. Το \lexiword{πρίσμα} βγαίνει από το αρχαίο «πρίζω», που σημαίνει πριονίζω: κάτι «πριονισμένο». Το \lexiword{παραλληλεπίπεδο} έχει «παράλληλα επίπεδα» για απέναντι έδρες. Ο \lexiword{κώνος} είναι το κουκουνάρι στα αρχαία ελληνικά. Ο \lexiword{κύλινδρος} βγαίνει από το «κυλίω». Η \lexiword{γενέτειρα} είναι αυτή που «γεννά» την επιφάνεια όταν περιστρέφεται. Το \lexiword{ανάπτυγμα} βγαίνει από το «αναπτύσσω», δηλαδή ξεδιπλώνω.
\end{lexi}

\begin{metafrasi}
Η «πλευρά» ενός κουτιού. & έδρα\\
Εκεί που συναντιούνται δύο έδρες. & ακμή\\
Εκεί που συναντιούνται οι ακμές (η «γωνία» του κουτιού). & κορυφή\\
Το «πάνω και το κάτω» ενός πρίσματος. & οι δύο βάσεις\\
Οι «τοίχοι» ενός πρίσματος. & παράπλευρες έδρες\\
Το κουτί ανοιγμένο και απλωμένο στο τραπέζι. & ανάπτυγμα\\
Η απόσταση ανάμεσα στις δύο βάσεις. & ύψος\\
\end{metafrasi}

\tmima{Η θεωρία}

\ypotmima{Α}{Έδρες, ακμές, κορυφές: το ορθό πρίσμα}

Ένα ζάρι είναι ένα στερεό. Έχει $6$ \textbf{έδρες} (οι επίπεδες επιφάνειές του, εδώ τετράγωνα), $12$ \textbf{ακμές} (τα ευθύγραμμα τμήματα όπου συναντιούνται δύο έδρες) και $8$ \textbf{κορυφές} (τα σημεία όπου συναντιούνται οι ακμές).

\begin{center}
\begin{tikzpicture}[scale=0.9]
\binCuboid{2}{2}{1}{0.8}
\node[font=\small] (ed) at (-1.3,2.9) {έδρα};
\draw[-{Stealth}] (ed)--(0.7,1.2);
\node[font=\small] (ak) at (4.9,0.4) {ακμή};
\draw[-{Stealth}] (ak)--(3.05,0.9);
\node[font=\small] (ko) at (4.9,3.2) {κορυφή};
\draw[-{Stealth}] (ko)--(3.1,2.85);
\end{tikzpicture}
\end{center}

Όταν σχεδιάζουμε ένα στερεό, οι ακμές που δεν φαίνονται από εμάς (γιατί είναι πίσω από κάποια έδρα) τις γράφουμε διακεκομμένες.

\begin{orismos}{ορθό πρίσμα}
\textbf{Ορθό πρίσμα} λέγεται το στερεό που έχει:\\
$\bullet$ δύο παράλληλες έδρες που είναι ίσα πολύγωνα και λέγονται \textbf{βάσεις} του πρίσματος, και\\
$\bullet$ όλες τις άλλες έδρες ορθογώνια παραλληλόγραμμα, που λέγονται \textbf{παράπλευρες έδρες}.\\
Η απόσταση των δύο βάσεων, που είναι ίση με το ύψος μιας παράπλευρης έδρας, λέγεται \textbf{ύψος} του πρίσματος.
\end{orismos}

Στη συνέχεια τα ορθά πρίσματα θα τα λέμε απλώς πρίσματα. Το όνομα του πρίσματος εξαρτάται από τη βάση του: αν η βάση είναι τρίγωνο, τετράπλευρο, πεντάγωνο κ.ο.κ., το πρίσμα λέγεται \textbf{τριγωνικό}, \textbf{τετραπλευρικό}, \textbf{πενταγωνικό} κ.ο.κ.

\begin{center}
\begin{tikzpicture}[scale=0.8]
\begin{scope}\binPrism{3}{1.3}{0.35}{2.2}{-90}
\node[font=\small,below] at (0,-0.7) {τριγωνικό};\end{scope}
\begin{scope}[shift={(4,0)}]\binPrism{5}{1.3}{0.35}{2.2}{-90}
\node[font=\small,below] at (0,-0.7) {πενταγωνικό};\end{scope}
\begin{scope}[shift={(8,0)}]\binPrism{6}{1.3}{0.35}{2.2}{-60}
\node[font=\small,below] at (0,-0.7) {εξαγωνικό};\end{scope}
\node[font=\small] (b1) at (-2,3.4) {βάση};\draw[-{Stealth}] (b1)--(-0.5,2.4);
\node[font=\small,align=center] (p1) at (-2.4,0.4) {παράπλευρη\\έδρα};\draw[-{Stealth}] (p1)--(-0.6,1.0);
\end{tikzpicture}
\end{center}

\begin{prosoxi}
Στο πρίσμα οι \emph{βάσεις} είναι οι δύο ίσες και παράλληλες έδρες, όχι οι έδρες «από κάτω». Αν ακουμπήσεις ένα πρίσμα στο τραπέζι από μια παράπλευρη έδρα, οι βάσεις του δεν αλλάζουν.
\end{prosoxi}

\tmimasmall{Το ανάπτυγμα}

Αν ξεδιπλώσουμε ένα πρίσμα και το απλώσουμε σε ένα επίπεδο, παίρνουμε το \textbf{ανάπτυγμά} του. Αποτελείται από τις παράπλευρες έδρες, που σχηματίζουν ένα ορθογώνιο με μήκος ίσο με την περίμετρο της βάσης και πλάτος ίσο με το ύψος του πρίσματος, και από τις δύο βάσεις.

\begin{center}
\begin{tikzpicture}[scale=0.62]
\foreach \x in {0,1.8,3.6}{\draw[fill=gblue!8,thick] (\x,0) rectangle (\x+1.8,3.6);}
\draw[fill=gblue!20,thick] (1.8,3.6)--(3.6,3.6)--(2.7,3.6+1.56)--cycle;
\draw[fill=gblue!20,thick] (1.8,0)--(3.6,0)--(2.7,-1.56)--cycle;
\node at (0.9,1.8) {έδρα};\node at (2.7,1.8) {έδρα};\node at (4.5,1.8) {έδρα};
\node at (2.7,4.1) {βάση};\node at (2.7,-0.8) {βάση};
\draw[<->] (0,-2.0)--(5.4,-2.0) node[midway,below,font=\small] {περίμετρος της βάσης};
\draw[<->] (-0.4,0)--(-0.4,3.6) node[midway,left,font=\small] {ύψος};
\end{tikzpicture}
\end{center}

\begin{kanonas}{ανάπτυγμα πρίσματος}
Το ανάπτυγμα ενός πρίσματος έχει ένα ορθογώνιο με διαστάσεις (περίμετρος βάσης) $\times$ (ύψος), και δύο ίσες βάσεις, μία από πάνω και μία από κάτω.
\end{kanonas}

\tmimasmall{Παραλληλεπίπεδο και κύβος}

\begin{orismos}{παραλληλεπίπεδο}
\textbf{Παραλληλεπίπεδο} λέγεται το πρίσμα που έχει βάσεις παραλληλόγραμμα. Αν οι βάσεις είναι ορθογώνια, όλες οι έδρες του είναι ορθογώνια και το λέμε \textbf{ορθογώνιο παραλληλεπίπεδο}. Αν όλες οι έδρες είναι ίσα τετράγωνα, το λέμε \textbf{κύβο}.
\end{orismos}

Το ορθογώνιο παραλληλεπίπεδο και ο κύβος έχουν $6$ έδρες, $12$ ακμές και $8$ κορυφές. Οι απέναντι έδρες είναι ίσες και οι $12$ ακμές χωρίζονται σε τρεις ομάδες των τεσσάρων ίσων ακμών: το μήκος, το πλάτος και το ύψος. Στον κύβο όλες οι ακμές είναι ίσες.

\begin{center}
\begin{tikzpicture}[scale=0.55]
\begin{scope}\binCuboid{2.4}{2.4}{1.2}{0.9}
\node[font=\small] at (1.8,-0.6) {κύβος};\end{scope}
\begin{scope}[shift={(6.2,0)}]\binCuboid{4}{1.6}{1.2}{0.9}
\node[font=\small] at (2.6,-0.6) {ορθογώνιο παραλληλεπίπεδο};\end{scope}
\end{tikzpicture}
\end{center}

Υπάρχουν $11$ διαφορετικά αναπτύγματα του κύβου: έξι τετράγωνα που διπλώνονται και σχηματίζουν κύβο, χωρίς να πέφτουν το ένα πάνω στο άλλο.

\begin{center}
\begin{tikzpicture}[scale=0.42]
\binNet{1/0,0/1,1/1,2/1,1/2,1/3}
\begin{scope}[shift={(6,0)}]\binNet{0/0,1/0,2/0,3/0,1/1,3/-1}\end{scope}
\begin{scope}[shift={(12,0)}]\binNet{0/1,1/1,1/0,2/0,2/-1,3/-1}\end{scope}
\end{tikzpicture}
\quad\quad
\begin{tikzpicture}[scale=0.4]
\draw[fill=gblue!12,thick] (0,0) rectangle (3,1.5);
\draw[fill=gblue!12,thick] (0,1.5) rectangle (3,3.5);
\draw[fill=gblue!12,thick] (0,3.5) rectangle (3,5);
\draw[fill=gblue!12,thick] (0,5) rectangle (3,6.5);
\draw[fill=gblue!12,thick] (-1.5,1.5) rectangle (0,3.5);
\draw[fill=gblue!12,thick] (3,1.5) rectangle (4.5,3.5);
\end{tikzpicture}
\end{center}
Τρία από τα αναπτύγματα του κύβου (αριστερά) και ένα ανάπτυγμα ορθογώνιου παραλληλεπιπέδου (δεξιά).

\begin{prosoxi}
Έξι τετράγωνα σε οποιαδήποτε διάταξη δεν δίνουν πάντα κύβο. Όπως ανακάλυψε η Αγγελική, αν τα βάλεις σε σειρά (ή αν δύο από αυτά θα έπεφταν στην ίδια έδρα), δεν κλείνει. Για να ελέγξεις ένα ανάπτυγμα, δίπλωσέ το με χαρτί.
\end{prosoxi}

\ypotmima{Β}{Η πυραμίδα}

\begin{orismos}{πυραμίδα}
\textbf{Πυραμίδα} λέγεται το στερεό που έχει μία έδρα πολύγωνο, τη \textbf{βάση}, και όλες τις άλλες έδρες τρίγωνα με κοινή κορυφή. Οι έδρες αυτές λέγονται \textbf{παράπλευρες έδρες}, και η κοινή τους κορυφή \textbf{κορυφή της πυραμίδας}. Το κάθετο ευθύγραμμο τμήμα από την κορυφή προς το επίπεδο της βάσης λέγεται \textbf{ύψος} της πυραμίδας.
\end{orismos}

Η πυραμίδα παίρνει όνομα από τη βάση της: \textbf{τριγωνική} (λέγεται και \textbf{τετράεδρο}), \textbf{τετραγωνική}, \textbf{πενταγωνική} κ.ο.κ.

\begin{center}
\begin{tikzpicture}[scale=0.85]
\begin{scope}\binPyr{4}{1.5}{0.35}{2.6}{-90}
\node[font=\small] (k) at (-2.2,2.9) {κορυφή};\draw[-{Stealth}] (k)--(-0.2,2.55);
\node[font=\small] (y) at (-2.4,1.2) {ύψος};\draw[-{Stealth}] (y)--(-0.08,1.2);
\draw[densely dotted] (0,0)--(0,2.6);
\node[font=\small] (b) at (2.3,-0.9) {βάση};\draw[-{Stealth}] (b)--(1.0,-0.2);
\node[font=\small,align=center] (p) at (2.4,2.0) {παράπλευρη\\έδρα};\draw[-{Stealth}] (p)--(0.6,0.9);
\end{scope}
\end{tikzpicture}
\end{center}

\begin{center}
\begin{tikzpicture}[scale=0.8]
\begin{scope}
\draw[dashed,gray] (0,0)--(3.2,0.8);
\draw[thick] (0,0)--(2.2,0)--(3.2,0.8)--(1.3,2.2)--cycle;
\draw[thick] (1.3,2.2)--(2.2,0);
\node[font=\small,below] at (1.6,-0.1) {τετράεδρο};
\end{scope}
\begin{scope}[shift={(5,0)}]\binPyr{4}{1.3}{0.35}{2.2}{-90}
\node[font=\small,below] at (0,-0.5) {τετραγωνική};\end{scope}
\begin{scope}[shift={(8.5,0)}]\binPyr{5}{1.3}{0.35}{2.2}{-90}
\node[font=\small,below] at (0,-0.5) {πενταγωνική};\end{scope}
\end{tikzpicture}
\end{center}

\begin{center}
\begin{tikzpicture}[scale=0.55]
\draw[fill=gblue!12,thick] (0,0) rectangle (2.4,2.4);
\draw[fill=gblue!20,thick] (0,2.4)--(2.4,2.4)--(1.2,5.4)--cycle;
\draw[fill=gblue!20,thick] (2.4,0)--(2.4,2.4)--(5.4,1.2)--cycle;
\draw[fill=gblue!20,thick] (0,0)--(2.4,0)--(1.2,-3)--cycle;
\draw[fill=gblue!20,thick] (0,0)--(0,2.4)--(-3,1.2)--cycle;
\begin{scope}[shift={(10,0)}]
\draw[fill=gblue!12,thick] (0,0)--(3.6,0)--(1.8,3.1)--cycle;
\draw[fill=gblue!20,thick] (0,0)--(3.6,0)--(1.8,-3.1)--cycle;
\draw[fill=gblue!20,thick] (0,0)--(1.8,3.1)--(-1.8,3.1)--cycle;
\draw[fill=gblue!20,thick] (3.6,0)--(1.8,3.1)--(5.4,3.1)--cycle;
\end{scope}
\end{tikzpicture}
\end{center}
Ανάπτυγμα τετραγωνικής πυραμίδας (αριστερά) και τετραέδρου (δεξιά).

Αν η βάση μιας πυραμίδας ή ενός πρίσματος έχει $\nu$ πλευρές, μπορούμε να μετρήσουμε τα στοιχεία τους με έναν κανόνα:

\begin{center}
\renewcommand{\arraystretch}{1.4}
\begin{tabular}{@{}lccc@{}}
\toprule
 & \textbf{Έδρες (Ε)} & \textbf{Κορυφές (Κ)} & \textbf{Ακμές (Α)}\\
\midrule
$\nu$-γωνική πυραμίδα & $\nu+1$ & $\nu+1$ & $2\nu$\\
$\nu$-γωνικό πρίσμα & $\nu+2$ & $2\nu$ & $3\nu$\\
\bottomrule
\end{tabular}
\end{center}

\begin{prosoxi}
Το ύψος της πυραμίδας είναι το κάθετο τμήμα προς τη βάση, και όχι το ύψος μιας παράπλευρης έδρας. Στην ίδια πυραμίδα το πρώτο είναι πάντα μικρότερο από το δεύτερο. Επίσης η πυραμίδα έχει μία μόνο βάση, ενώ το πρίσμα δύο.
\end{prosoxi}

\ypotmima{Γ}{Κύλινδρος, κώνος και σφαίρα}

Στη ζωή συναντάμε και στερεά με καμπύλες επιφάνειες: κεριά, κουτάκια, χωνάκια, μπάλες. Μπορούμε να τα δούμε ως σχήματα που παράγονται όταν ένα επίπεδο σχήμα \emph{περιστρέφεται} γύρω από μια ευθεία.

\begin{orismos}{κύλινδρος}
\textbf{Κύλινδρος} λέγεται το στερεό που παράγεται από την περιστροφή ενός ορθογωνίου γύρω από μία πλευρά του. Η πλευρά που περιστρέφεται λέγεται \textbf{γενέτειρα}. Ο κύλινδρος έχει δύο ίσους και παράλληλους κυκλικούς δίσκους, τις \textbf{βάσεις}, και μία \textbf{παράπλευρη (κυρτή) επιφάνεια}. Η απόσταση των βάσεων είναι το \textbf{ύψος} του.
\end{orismos}

Αν ξεδιπλώσουμε τον κύλινδρο, η παράπλευρη επιφάνεια γίνεται ορθογώνιο. Το μήκος του ορθογωνίου ισούται με την περιφέρεια της βάσης, δηλαδή $2\pi\rho\approx6{,}28\rho$ (από τις προηγούμενες τάξεις γνωρίζουμε ότι η περιφέρεια κύκλου ακτίνας $\rho$ είναι $2\pi\rho$, με $\pi\approx3{,}14$), και το πλάτος του είναι το ύψος του κυλίνδρου.

\begin{center}
\begin{tikzpicture}[scale=0.8]
\begin{scope}\binCyl{1}{0.3}{2.2}
\node[font=\small,left] at (-1.1,1.1) {ύψος};
\node[font=\small] at (0,-0.65) {κύλινδρος};
\end{scope}
\begin{scope}[shift={(5,-1.3)}]
\draw[thick,fill=gblue!8] (0,0) rectangle (3.8,2.2);
\draw[thick,fill=gblue!18] (1.2,2.2+0.6) circle (0.6);
\draw[thick,fill=gblue!18] (2.6,-0.6) circle (0.6);
\node[font=\small] at (1.9,1.1) {παράπλευρη επιφάνεια};
\draw[<->] (0,-1.5)--(3.8,-1.5) node[midway,below,font=\small] {περιφέρεια της βάσης};
\end{scope}
\end{tikzpicture}
\end{center}

\begin{orismos}{κώνος}
\textbf{Κώνος} λέγεται το στερεό που παράγεται από την περιστροφή ενός ορθογωνίου τριγώνου $\mathrm{ΟΑΒ}$ (με ορθή γωνία στο $\mathrm{Ο}$) γύρω από μία κάθετη πλευρά του, την $\mathrm{ΟΑ}$. Η $\mathrm{ΟΑ}$ είναι το \textbf{ύψος} του κώνου. Ο κώνος έχει μία \textbf{βάση}, έναν κυκλικό δίσκο με κέντρο $\mathrm{Ο}$ και \textbf{ακτίνα} $\mathrm{ΟΒ}$, και μία \textbf{παράπλευρη επιφάνεια}. Η πλευρά $\mathrm{ΑΒ}$ είναι η \textbf{γενέτειρα}.
\end{orismos}

Η παράπλευρη επιφάνεια του κώνου ξεδιπλώνεται σε έναν \textbf{κυκλικό τομέα} με ακτίνα τη γενέτειρα $\mathrm{ΑΒ}$. Το μήκος του τόξου του ισούται με την περιφέρεια της βάσης, και η γωνία του τομέα ($\mu^\circ$) βρίσκεται από τη σχέση $\mu=\dfrac{360\cdot\rho}{\mathrm{ΑΒ}}$, όπου $\rho$ η ακτίνα της βάσης.

\begin{center}
\begin{tikzpicture}[scale=0.8]
\begin{scope}
\binCone{1}{0.3}{2.4}
\node[font=\small,anchor=east] at (-0.1,1.2) {ύψος};
\node[font=\small,anchor=west] at (0.55,1.5) {γενέτειρα};
\draw[dashed,gray] (0,0)--(0,2.4);
\node[font=\small] at (0,-0.65) {κώνος};
\end{scope}
\begin{scope}[shift={(6,2.4)}]
\draw[thick,fill=gblue!8] (0,0)--(198:2.5) arc[start angle=198,end angle=342,radius=2.5]--cycle;
\draw[thick,fill=gblue!18] (0,-3.5) circle (1);
\node[font=\small] at (0,-1.2) {$\mu$};
\node[font=\small,right] at (1.5,-1) {γενέτειρα};
\end{scope}
\end{tikzpicture}
\end{center}

\begin{orismos}{σφαίρα}
\textbf{Σφαίρα} λέγεται το στερεό που παράγεται από την περιστροφή ενός κύκλου με κέντρο $\mathrm{Ο}$ και ακτίνα $\rho$ γύρω από μία διάμετρό του. Το $\mathrm{Ο}$ λέγεται \textbf{κέντρο} και το $\rho$ \textbf{ακτίνα} της σφαίρας.
\end{orismos}

\begin{center}
\begin{tikzpicture}[scale=0.9]
\binSph{1.3}
\fill (0,0) circle (1.5pt) node[below left,font=\small] {$\mathrm{Ο}$};
\draw (0,0)--(1.3,0) node[midway,above,font=\small] {$\rho$};
\end{tikzpicture}
\end{center}

Η σφαίρα δεν έχει ούτε έδρες ούτε ακμές ούτε κορυφές, και δεν έχει ανάπτυγμα: δεν μπορείς να ξεδιπλώσεις μια μπάλα σε επίπεδο χωρίς να τη σχίσεις ή να την τραβήξεις. Ο κύλινδρος και ο κώνος δεν είναι πολύεδρα, γιατί έχουν καμπύλη επιφάνεια.

\begin{prosoxi}
Ο κύκλος είναι επίπεδο σχήμα και η σφαίρα στερεό: η μπάλα δεν είναι «κύκλος». Επίσης στον κώνο το ύψος και η γενέτειρα είναι διαφορετικά τμήματα: το ύψος είναι κάθετο στη βάση, ενώ η γενέτειρα είναι πάντα μεγαλύτερη από αυτό.
\end{prosoxi}

\ypotmima{Δ}{Όψεις και σχεδίαση στερεών}

Για να περιγράψουμε ένα στερεό σε χαρτί, το «κοιτάμε» από διαφορετικές κατευθύνσεις και σχεδιάζουμε αυτό που βλέπουμε. Κάθε τέτοιο σχέδιο λέγεται \textbf{όψη}: η \emph{μπροστινή} όψη (από μπροστά), η \emph{πάνω} όψη (από ψηλά) και η \emph{πλάγια} όψη (από το πλάι). Οι όψεις ενός κουτιού είναι ορθογώνια, της πυραμίδας τρίγωνα και τετράγωνο με τις διαγωνίους, του κυλίνδρου ορθογώνιο και κύκλος.

\begin{center}
\begin{tikzpicture}[scale=0.55]
\begin{scope}\binCuboid{3}{2}{1.5}{1}
\node[font=\small] at (2.2,-0.7) {στερεό};\end{scope}
\begin{scope}[shift={(6,0)}]
\draw[thick,fill=gblue!8] (0,0) rectangle (3,2);\node[font=\small] at (1.5,-0.7) {μπροστινή};\end{scope}
\begin{scope}[shift={(6,3.5)}]
\draw[thick,fill=gblue!8] (0,0) rectangle (3,1.5);\node[font=\small] at (1.5,2.1) {πάνω};\end{scope}
\begin{scope}[shift={(10.5,0)}]
\draw[thick,fill=gblue!8] (0,0) rectangle (1.5,2);\node[font=\small] at (0.75,-0.7) {πλάγια};\end{scope}
\begin{scope}[shift={(15,0)}]
\draw[thick,fill=gblue!8] (0,0)--(3,0)--(1.5,2.8)--cycle;\node[font=\small] at (1.5,-0.7) {πυραμίδα, μπροστά};
\draw[thick,fill=gblue!8] (0,3.5) rectangle (3,6.5);\draw (0,3.5)--(3,6.5) (0,6.5)--(3,3.5);\node[font=\small] at (1.5,7.1) {πάνω};\end{scope}
\end{tikzpicture}
\end{center}

Για να σχεδιάσουμε το ίδιο το στερεό χρησιμοποιούμε \textbf{ισομετρικό χαρτί}, που έχει τελείες σε κανονική τριγωνική διάταξη, ή τετραγωνισμένο χαρτί και διακεκομμένες γραμμές για τις κρυμμένες ακμές. Μπορούμε επίσης να χρησιμοποιήσουμε ψηφιακά εργαλεία γεωμετρίας (π.χ. GeoGebra) ή απλά ελεύθερη σχεδίαση.

\tmima{Λυμένα παραδείγματα}

\paradeigma{Ονομάζω τα στερεά}
\tixero Τρία αντικείμενα: ένα κεράκι τούρτας, ένα κουτί με σοκολάτες σε σχήμα τριγώνου και ένας κώνος τροχαίας.\par
\tithelo Το είδος του στερεού που θυμίζει καθένα.\par
\possindeo Το κεράκι έχει δύο ίσες κυκλικές βάσεις: κύλινδρος. Το κουτί σοκολάτας έχει δύο τριγωνικές βάσεις και τρία ορθογώνια: τριγωνικό πρίσμα. Ο κώνος τροχαίας έχει μία κυκλική βάση και μία κορυφή: κώνος.\par
\apantisi Κύλινδρος, τριγωνικό πρίσμα, κώνος.\par
\elegxos Στο πρίσμα οι παράπλευρες έδρες είναι ορθογώνια, ενώ στην πυραμίδα τρίγωνα.

\paradeigma{Μετρώ έδρες, ακμές και κορυφές}
\tixero Ένα εξαγωνικό πρίσμα (όπως ένα μολύβι) και μια πυραμίδα με βάση εννιάγωνο.\par
\tithelo Έδρες, κορυφές και ακμές καθενός.\par
\possindeo Στο εξαγωνικό πρίσμα: $\nu=6$, άρα $6+2=8$ έδρες (δύο βάσεις και έξι ορθογώνια), $2\cdot6=12$ κορυφές και $3\cdot6=18$ ακμές. Στην εννεαγωνική πυραμίδα $\nu=9$: $9+1=10$ έδρες, $9+1=10$ κορυφές και $2\cdot9=18$ ακμές.\par
\apantisi Πρίσμα: $8$ έδρες, $12$ κορυφές, $18$ ακμές. Πυραμίδα: $10$ έδρες, $10$ κορυφές, $18$ ακμές.\par
\elegxos Στο πρίσμα: $12$ ακμές των βάσεων ($6+6$) και $6$ παράπλευρες, σύνολο $18$.

\paradeigma{Ποιο ανάπτυγμα κλείνει σε κύβο;}
\tixero Δύο διατάξεις από έξι τετράγωνα: οι έξι σε μια σειρά, και ο σταυρός από την αρχή της παραγράφου.\par
\tithelo Ποια διπλώνεται σε κύβο.\par
\possindeo Στη σειρά, όταν διπλώνω, το πρώτο και το πέμπτο τετράγωνο πέφτουν στην ίδια έδρα. Στον σταυρό, τέσσερα τετράγωνα σχηματίζουν τα πλάγια του κύβου και ένα από πάνω και ένα από κάτω καλύπτουν τις βάσεις.\par
\apantisi Μόνο ο σταυρός.\par
\elegxos Κόβω χαρτόνι και δίπλωνω: ο σταυρός κλείνει χωρίς κενά.

\paradeigma{Ανάπτυγμα τριγωνικού πρίσματος}
\tixero Πρίσμα με ισόπλευρη βάση πλευράς $3$ cm και ύψος $8$ cm.\par
\tithelo Τις διαστάσεις του ορθογωνίου της παράπλευρης επιφάνειας και τα υπόλοιπα κομμάτια του αναπτύγματος.\par
\possindeo Η περίμετρος της βάσης είναι $3\cdot3=9$ cm. Το ορθογώνιο έχει διαστάσεις $9\times8$ cm και αποτελείται από τρεις παράπλευρες έδρες $3\times8$ cm. Επιπλέον χρειάζονται δύο ισόπλευρα τρίγωνα πλευράς $3$ cm.\par
\apantisi Ορθογώνιο $9\times8$ cm και δύο ισόπλευρα τρίγωνα με πλευρά $3$ cm.\par
\elegxos $3$ έδρες $+$ $2$ βάσεις $=5$ έδρες, όσες και του τριγωνικού πρίσματος ($3+2$).

\paradeigma{Ανάπτυγμα κυλίνδρου}
\tixero Κύλινδρος με ακτίνα βάσης $\rho=2$ cm και ύψος $5$ cm.\par
\tithelo Τις διαστάσεις του αναπτύγματος.\par
\possindeo Η περιφέρεια της βάσης είναι $2\cdot3{,}14\cdot2=12{,}56\approx12{,}6$ cm. Το ορθογώνιο έχει διαστάσεις $12{,}6\times5$ cm και υπάρχουν δύο κύκλοι ακτίνας $2$ cm.\par
\apantisi Ορθογώνιο $12{,}6\times5$ cm και δύο κύκλοι με ακτίνα $2$ cm.\par
\elegxos Αν τυλίξουμε το ορθογώνιο, η μια πλευρά του $12{,}6$ cm κάνει ακριβώς έναν γύρο γύρω από τον κύκλο.

\paradeigma{Ανάπτυγμα κώνου}
\tixero Κώνος με ακτίνα βάσης $\rho=2$ cm και γενέτειρα $5$ cm.\par
\tithelo Το ανάπτυγμά του.\par
\possindeo Η παράπλευρη επιφάνεια είναι κυκλικός τομέας με ακτίνα $5$ cm και γωνία $\mu=\dfrac{360\cdot2}{5}=144^\circ$. Η βάση είναι κύκλος ακτίνας $2$ cm.\par
\apantisi Τομέας ακτίνας $5$ cm και γωνίας $144^\circ$ και κύκλος ακτίνας $2$ cm.\par
\elegxos Το μήκος του τόξου του τομέα είναι τα $\frac{144}{360}=\frac{2}{5}$ του κύκλου ακτίνας $5$, δηλαδή $\frac{2}{5}\cdot2\pi\cdot5=2\pi\cdot2$, όσο η περιφέρεια της βάσης.

\paradeigma{Περιστροφές}
\tixero Ένα ορθογώνιο με πλευρές $3$ cm και $5$ cm, και ένα ορθογώνιο τρίγωνο με κάθετες πλευρές $3$ cm και $4$ cm περιστρέφονται γύρω από την πλευρά των $5$ cm και των $4$ cm αντίστοιχα.\par
\tithelo Τα στερεά που προκύπτουν.\par
\possindeo Η περιστροφή ορθογωνίου δίνει κύλινδρο: με ύψος $5$ cm (η πλευρά-άξονας) και ακτίνα βάσης $3$ cm. Η περιστροφή ορθογωνίου τριγώνου γύρω από κάθετη πλευρά δίνει κώνο: με ύψος $4$ cm και ακτίνα βάσης $3$ cm.\par
\apantisi Κύλινδρος ($\rho=3$, ύψος $5$) και κώνος ($\rho=3$, ύψος $4$).\par
\elegxos Η πλευρά που περιστρέφεται (γενέτειρα) στον κύλινδρο είναι παράλληλη στον άξονα, ενώ στον κώνο είναι η υποτείνουσα.

\paradeigma{Οι όψεις ενός στερεού}
\tixero Μια τετραγωνική πυραμίδα με βάση $4$ cm και ύψος $5$ cm.\par
\tithelo Τις όψεις της από πάνω και από μπροστά.\par
\possindeo Από πάνω βλέπω το τετράγωνο της βάσης και τις τέσσερις ακμές από την κορυφή προς τις γωνίες του, δηλαδή τις δύο διαγωνίους. Από μπροστά βλέπω ένα τρίγωνο με βάση $4$ cm και ύψος $5$ cm.\par
\apantisi Πάνω: τετράγωνο $4\times4$ με διαγωνίους. Μπροστά: τρίγωνο με βάση $4$ cm και ύψος $5$ cm.\par
\elegxos Στην όψη από πάνω η κορυφή βρίσκεται στο κέντρο του τετραγώνου.

\tmima{Ασκήσεις}

\zoni{Το κατάλαβα}{ονόματα, στοιχεία, αναπτύγματα}

\ex Συμπλήρωσε τα κενά. α) Οι βάσεις ενός πρίσματος είναι δύο \dots\ πολύγωνα. \quad β) Η απόσταση των δύο βάσεων ονομάζεται \dots\ του πρίσματος. \quad γ) Οι παράπλευρες έδρες ενός ορθού πρίσματος είναι \dots. \quad δ) Η κοινή κορυφή των τριγώνων μιας πυραμίδας λέγεται \dots\ της πυραμίδας. \quad ε) Ο κύλινδρος έχει ως βάσεις δύο ίσους \dots\ δίσκους. \quad στ) Ο κώνος έχει μία \dots\ που είναι ένας κυκλικός δίσκος. \quad ζ) Κατά την περιστροφή ενός κύκλου γύρω από μια διάμετρό του παράγεται \dots.

\ex Πες ποιο στερεό θυμίζει κάθε αντικείμενο: α) ζάρι \quad β) κουτί παπουτσιών \quad γ) κουτάκι αναψυκτικού \quad δ) χωνάκι παγωτού \quad ε) μπάλα ποδοσφαίρου \quad στ) σκηνή με τριγωνικές πλευρές \quad ζ) καπέλο πάρτι \quad η) αιγυπτιακή πυραμίδα \quad θ) εξαγωνικό μολύβι.

\ex Ποιο στερεό περιγράφεται; α) Έχει $6$ έδρες, όλες τετράγωνα. \quad β) Έχει δύο ίσους κυκλικούς δίσκους και μία κυρτή επιφάνεια. \quad γ) Έχει μία κυκλική βάση και μία κορυφή. \quad δ) Έχει $5$ έδρες: μία τετράγωνη και τέσσερις τριγωνικές. \quad ε) Δεν έχει ακμές ούτε κορυφές και όλα τα σημεία της απέχουν το ίδιο από το κέντρο. \quad στ) Έχει $5$ έδρες: δύο τριγωνικές και τρεις ορθογώνιες.

\ex Σωστό ή λάθος; α) Ο κύβος έχει $12$ ακμές. \quad β) Κάθε πρίσμα έχει δύο βάσεις. \quad γ) Η πυραμίδα έχει δύο βάσεις. \quad δ) Οι παράπλευρες έδρες της πυραμίδας είναι ορθογώνια. \quad ε) Ο κώνος έχει δύο βάσεις. \quad στ) Η σφαίρα έχει μία κορυφή. \quad ζ) Το ανάπτυγμα του κυλίνδρου αποτελείται από ένα ορθογώνιο και δύο κύκλους.

\ex Ένα κουτί σε σχήμα ορθογωνίου παραλληλεπιπέδου έχει διαστάσεις $6$, $4$ και $3$ cm. α) Πόσες έδρες του είναι ίσες με την έδρα $6\times4$ και πόσες είναι οι άλλες έδρες και τι διαστάσεις έχουν; β) Πόσες ακμές έχει από κάθε μήκος;

\ex Ονόμασε το στερεό και πες πόσες παράπλευρες έδρες έχει: α) πρίσμα με βάση πεντάγωνο \quad β) πυραμίδα με βάση τετράπλευρο \quad γ) πρίσμα με βάση εξάγωνο \quad δ) πυραμίδα με βάση τρίγωνο.

\ex \textbf{Αναπτύγματα κύβου.} Ποια από τις παρακάτω διατάξεις έξι τετραγώνων διπλώνονται σε κύβο;
\begin{center}
\begin{tikzpicture}[scale=0.42]
\node[font=\small] at (1.5,-1.4) {α};
\begin{scope}\binNet{1/0,1/1,1/2,1/3,0/1,2/1}\end{scope}
\begin{scope}[shift={(7,0)}]\node[font=\small] at (2,-1.4) {β};\binNet{0/0,1/0,2/0,3/0,0/1,3/1}\end{scope}
\begin{scope}[shift={(14,0)}]\node[font=\small] at (2,-1.4) {γ};\binNet{0/0,1/0,2/0,3/0,1/1,2/-1}\end{scope}
\begin{scope}[shift={(0,-8)}]\node[font=\small] at (2.5,-1.4) {δ};\binNet{0/0,1/0,2/0,3/0,4/0,2/1}\end{scope}
\begin{scope}[shift={(7,-8)}]\node[font=\small] at (1.5,-1.4) {ε};\binNet{0/0,1/0,1/1,2/1,2/2,3/2}\end{scope}
\begin{scope}[shift={(14,-8)}]\node[font=\small] at (1.5,-1.4) {στ};\binNet{0/0,1/0,2/0,0/1,1/1,2/1}\end{scope}
\end{tikzpicture}
\end{center}

\ex \textbf{Ανάπτυγμα και στερεό.} Αντιστοίχισε κάθε ανάπτυγμα με το στερεό που σχηματίζει: τριγωνικό πρίσμα, τετραγωνική πυραμίδα, κύλινδρος, κώνος.
\begin{center}
\begin{tikzpicture}[scale=0.34]
\node[font=\small] at (2.5,-2.4) {Α};
\draw[fill=gblue!8,thick] (0,0) rectangle (5,2.4);\draw[fill=gblue!18,thick] (1.2,2.4+0.8) circle (0.8);\draw[fill=gblue!18,thick] (3.8,-0.8) circle (0.8);
\begin{scope}[shift={(8,0)}]\node[font=\small] at (2.4,-2.4) {Β};
\draw[fill=gblue!8,thick] (0,0) rectangle (2.4,2.4);
\draw[fill=gblue!18,thick] (0,2.4)--(2.4,2.4)--(1.2,4.8)--cycle;
\draw[fill=gblue!18,thick] (2.4,0)--(2.4,2.4)--(4.8,1.2)--cycle;
\draw[fill=gblue!18,thick] (0,0)--(2.4,0)--(1.2,-2.4)--cycle;
\draw[fill=gblue!18,thick] (0,0)--(0,2.4)--(-2.4,1.2)--cycle;\end{scope}
\begin{scope}[shift={(19,5)}]\node[font=\small] at (0,-6.7) {Γ};
\draw[thick,fill=gblue!8] (0,0)--(192:3) arc[start angle=192,end angle=348,radius=3]--cycle;
\draw[thick,fill=gblue!18] (0,-4.4) circle (1.4);\end{scope}
\begin{scope}[shift={(26,0)}]\node[font=\small] at (3,-2.4) {Δ};
\foreach \x in {0,2,4}{\draw[fill=gblue!8,thick] (\x,0) rectangle (\x+2,4);}
\draw[fill=gblue!18,thick] (2,4)--(4,4)--(3,5.7)--cycle;
\draw[fill=gblue!18,thick] (2,0)--(4,0)--(3,-1.7)--cycle;\end{scope}
\end{tikzpicture}
\end{center}

\ex \textbf{Τα στοιχεία της πυραμίδας.} Στο σχήμα φαίνεται μια τετραγωνική πυραμίδα. Πες ποια από τα στοιχεία \emph{κορυφή της πυραμίδας, βάση, ύψος, παράπλευρη έδρα} δείχνουν τα γράμματα $\alpha$, $\beta$, $\gamma$, $\delta$.
\begin{center}
\begin{tikzpicture}[scale=0.85]
\binPyr{4}{1.5}{0.35}{2.6}{-90}
\draw[densely dotted] (0,0)--(0,2.6);
\node (a) at (-2,3) {$\alpha$};\draw[-{Stealth}] (a)--(-0.15,2.6);
\node (b) at (2.4,-0.9) {$\beta$};\draw[-{Stealth}] (b)--(1.0,-0.2);
\node (c) at (-1.5,1.0) {$\gamma$};\draw[-{Stealth}] (c)--(-0.08,1.0);
\node (d) at (2.4,1.8) {$\delta$};\draw[-{Stealth}] (d)--(0.6,0.9);
\end{tikzpicture}
\end{center}

\zoni{Το χρησιμοποιώ}{πίνακες, αναπτύγματα, όψεις}

\ex Συμπλήρωσε τον πίνακα για πυραμίδες (το $\nu$ είναι το πλήθος των πλευρών της βάσης):
\begin{center}
\renewcommand{\arraystretch}{1.3}
\begin{tabular}{@{}lcccc@{}}
\toprule
\textbf{Πυραμίδα} & \textbf{Πλευρές βάσης} & \textbf{Έδρες (Ε)} & \textbf{Κορυφές (Κ)} & \textbf{Ακμές (Α)}\\
\midrule
Τετραγωνική & $4$ & $5$ & $5$ & $8$\\
Εξαγωνική & $6$ & & & \\
Οκταγωνική & $8$ & & & \\
Δωδεκαγωνική & $12$ & & & \\
$\nu$-γωνική & $\nu$ & & & \\
\bottomrule
\end{tabular}
\end{center}
Βρες για κάθε γραμμή το $\mathrm{Κ}+\mathrm{Ε}$ και το $\mathrm{Α}$. Τι παρατηρείς;

\ex Συμπλήρωσε τον πίνακα για πρίσματα:
\begin{center}
\renewcommand{\arraystretch}{1.3}
\begin{tabular}{@{}lcccc@{}}
\toprule
\textbf{Πρίσμα} & \textbf{Πλευρές βάσης} & \textbf{Έδρες (Ε)} & \textbf{Κορυφές (Κ)} & \textbf{Ακμές (Α)}\\
\midrule
Κύβος & $4$ & $6$ & $8$ & $12$\\
Τριγωνικό & $3$ & & & \\
Πενταγωνικό & $5$ & & & \\
Επταγωνικό & $7$ & & & \\
Δωδεκαγωνικό & $12$ & & & \\
$\nu$-γωνικό & $\nu$ & & & \\
\bottomrule
\end{tabular}
\end{center}
Έλεγξε ότι και εδώ ισχύει η σχέση που βρήκες στην προηγούμενη άσκηση.

\ex \textbf{Η κορδέλα.} Ο Στέλιος θέλει να τυλίξει κορδέλα κατά μήκος όλων των ακμών ενός κουτιού παπουτσιών με διαστάσεις $30$, $20$ και $12$ cm. Πόση κορδέλα χρειάζεται; Πόση θα χρειαζόταν για έναν κύβο ακμής $7$ cm;

\ex \textbf{Το πρίσμα της Φανής.} Ένα τριγωνικό πρίσμα έχει ισόπλευρη βάση πλευράς $4$ cm και ύψος $10$ cm. α) Πόσες και τι έδρες έχει; β) Ποιες διαστάσεις έχει το ορθογώνιο της παράπλευρης επιφάνειας του αναπτύγματός του; γ) Σχεδίασε το ανάπτυγμα σε τετραγωνισμένο χαρτί.

\ex \textbf{Ο κύλινδρος.} Κατασκεύασε τις όψεις και το ανάπτυγμα κυλίνδρου με ακτίνα βάσης $\rho=2$ cm, ύψος $6$ cm και περιφέρεια βάσης $12{,}6$ cm. Ποιες διαστάσεις έχει το ορθογώνιο του αναπτύγματος;

\ex \textbf{Ο κώνος.} Κατασκεύασε το ανάπτυγμα κώνου με ακτίνα βάσης $\rho=2$ cm και γενέτειρα $5$ cm. Η γωνία του κυκλικού τομέα είναι $\mu=144^\circ$.

\ex \textbf{Οι όψεις.} Ένα ορθογώνιο παραλληλεπίπεδο έχει διαστάσεις (μήκος, πλάτος, ύψος) $4$, $3$ και $2$ cm. Σχεδίασε τις τρεις όψεις του. Σχεδίασε και τις όψεις (από μπροστά και από πάνω) μιας τετραγωνικής πυραμίδας με βάση $4$ cm και ύψος $5$ cm.

\ex \textbf{Περιστροφές και σφαίρα.} Ποιο στερεό παράγεται και με ποια στοιχεία όταν περιστραφεί: α) ορθογώνιο $3\times5$ cm γύρω από την πλευρά των $5$ cm \quad β) ορθογώνιο τρίγωνο με κάθετες πλευρές $3$ και $4$ cm και υποτείνουσα $5$ cm γύρω από την κάθετη πλευρά των $4$ cm \quad γ) ημικύκλιο ακτίνας $5$ cm γύρω από τη διάμετρό του; Στο (γ), ποια είναι η μεγαλύτερη απόσταση δύο σημείων του στερεού; (Ο Νίκος κρατά μια μπάλα του μπάσκετ με διάμετρο $24$ cm. Ποια είναι η ακτίνα της;)

\zoni{Το σκέφτομαι}{ανάποδη σκέψη, ταξινόμηση, λάθη}

\ex \textbf{Από το πλήθος στο στερεό.} α) Ποιο πρίσμα έχει $15$ ακμές; \quad β) Ποια πυραμίδα έχει $12$ ακμές; \quad γ) Ποιο πρίσμα έχει $14$ κορυφές; \quad δ) Υπάρχει πρίσμα με $9$ κορυφές; Δικαιολόγησε.

\ex Ο Βασίλης λέει: «Ο κύβος είναι ορθογώνιο παραλληλεπίπεδο». Η Ρήνη λέει: «Και το ορθογώνιο παραλληλεπίπεδο είναι κύβος». Ποιος έχει δίκιο; Έλεγξε επίσης αν: α) κάθε ορθογώνιο παραλληλεπίπεδο είναι πρίσμα, β) κάθε πρίσμα είναι παραλληλεπίπεδο.

\ex Η Ζωή κόβει ανάπτυγμα κυλίνδρου με βάσεις ακτίνας $3$ cm και ορθογώνιο με μήκος $6$ cm (όσο η διάμετρος της βάσης). Θα κλείσει ο κύλινδρος; Πόσο μήκος χρειάζεται το ορθογώνιο;

\ex Στο σχήμα φαίνεται ένα ανάπτυγμα κύβου με έξι έδρες που τις ονομάσαμε $\mathrm{Α},\mathrm{Β},\Gamma,\Delta,\mathrm{Ε},\mathrm{Ζ}$. Ποιες έδρες θα βρεθούν απέναντι όταν διπλωθεί ο κύβος; Ποιες είναι δίπλα στην $\mathrm{Β}$;
\begin{center}
\begin{tikzpicture}[scale=0.55]
\binNet{0/0,1/0,2/0,3/0,1/1,2/-1}
\node at (0.5,0.5) {$\mathrm{Α}$};\node at (1.5,0.5) {$\mathrm{Β}$};\node at (2.5,0.5) {$\Gamma$};\node at (3.5,0.5) {$\Delta$};\node at (1.5,1.5) {$\mathrm{Ε}$};\node at (2.5,-0.5) {$\mathrm{Ζ}$};
\end{tikzpicture}
\end{center}

\ex \textbf{Ισομετρικό χαρτί.} Στο σχήμα ο κύβος ακμής $1$ έχει σχεδιαστεί σε ισομετρικό χαρτί. Σχεδίασε με τον ίδιο τρόπο ένα ορθογώνιο παραλληλεπίπεδο με διαστάσεις $2\times1\times1$ και μετά τις τρεις όψεις του.
\begin{center}
\begin{tikzpicture}[scale=0.8]
\foreach \i in {-3,...,3}{\foreach \j in {0,...,4}{
 \fill[gray] ({\i*0.866+(mod(\j,2))*0.433},{\j*0.5}) circle (0.7pt);}}
\draw[thick,fill=gblue!8] (0,0)--(0.866,0.5)--(0.866,1.5)--(0,1)--cycle;
\draw[thick,fill=gblue!14] (0,0)--(-0.866,0.5)--(-0.866,1.5)--(0,1)--cycle;
\draw[thick,fill=gblue!22] (0,1)--(0.866,1.5)--(0,2)--(-0.866,1.5)--cycle;
\end{tikzpicture}
\end{center}

\zoni{Για το σπίτι}{στερεά γύρω σου}

\ex Βρες στο σπίτι σου έξι αντικείμενα που έχουν σχήμα διαφορετικού στερεού (π.χ. κουτί δημητριακών, κονσέρβα, μπάλα, κώνος από πάρτι). Ζωγράφισέ τα και γράψε το όνομα του στερεού, και πόσες έδρες, ακμές και κορυφές έχει (όπου ισχύει).

\ex Φτιάξε από χαρτόνι ένα ανάπτυγμα κύβου με ακμή $5$ cm, κόψε το και δίπλωσέ το. Μετά φτιάξε και το ανάπτυγμα μιας τετραγωνικής πυραμίδας με βάση $4$ cm και παράπλευρες έδρες ισοσκελή τρίγωνα με ύψος $6$ cm.
